\section{实战推演：三类 Agent 项目的训练与交付路径}

前面完成了来源讲解和机制拆分，本节用三类项目检验这些原则能否落地。三条路径覆盖只读知识任务、有副作用的流程自动化和长时开放研究任务；比较重点不是选哪个模型，而是任务可验证性、工具权限、失败恢复和单位成功成本如何改变训练方案。

\subsection{路径 A：内部知识问答型 Agent}
这是最常见也是最容易被低估难度的一类。表面上像传统 RAG 问答，实际上只要涉及多轮澄清、跨系统检索与动作执行，就已经进入 Agent 设计范畴。基于本课方法，可采用 ``轻量 SFT + 任务化闭集评估 + 受限工具调用'' 的策略启动。

\begin{knowledgebox}{路径 A 的落地节奏}
\begin{enumerate}
    \item 第 1 周建立任务分解与工具白名单，只允许只读操作。
    \item 第 2 周补齐 close-ended evaluation，覆盖高频 FAQ 与关键业务术语。
    \item 第 3 周引入有限 open-ended 任务，跟踪失败恢复率与越权率。
    \item 第 4 周再扩充上下文窗口和工具集合，避免早期过度复杂化。
\end{enumerate}
\end{knowledgebox}

这种路径的成功关键不是追求最强模型，而是快速建立 ``可量化反馈''。若没有评估闭环，团队很容易在提示词和模板工程里反复迭代，却无法证明任务完成率是否真正改善。

\subsection{路径 B：流程自动化型 Agent}
路径 A 主要验证检索与回答，路径 B 则进入会改变业务状态的动作空间。写入工单、发起审批或修改 CRM 都要求幂等性、事务边界和回滚机制；训练数据也必须记录动作前后状态，不能只保存最终文本回复。
此类项目通常连接工单、审批、CRM、BI 等业务系统，目标是减少人工流程跳转。它对工具调用准确性和异常恢复能力要求高，训练策略应更偏向 Agent RL 与策略约束。

\begin{importantbox}{路径 B 的主要风险}
流程型 Agent 最怕 ``看似完成，实则破坏状态一致性''。例如自动填单成功率高，但字段语义错配；自动发起流程成功率高，但审批链漏节点。这类问题在仅看闭集准确率时很难暴露。
\end{importantbox}

\begin{table}[H]
\centering
\begin{tabular}{p{2.6cm}p{4.2cm}p{4.8cm}}
\toprule
\textbf{阶段} & \textbf{目标} & \textbf{验证重点} \\
\midrule
设计期 & 定义动作空间与权限边界 & 是否存在危险动作未隔离 \\
训练期 & 提升多步流程完成率 & 失败后是否能回滚并重计划 \\
试运行 & 降低人工介入频次 & 自动化收益是否覆盖推理成本 \\
扩张期 & 横向接入更多流程 & 新流程接入后旧流程是否退化 \\
\bottomrule
\end{tabular}
\caption{流程自动化型 Agent 的分阶段交付要点}
\end{table}

\subsection{路径 C：研究助理型 Agent}
流程自动化通常有明确终局状态，研究助理却面对多解、长证据链和动态工具。它需要同时评估检索覆盖、引用真实性、实验可复现性和结论边界；任何一个漂亮的最终报告，都必须能回溯到实际检索与执行记录。
研究助理型 Agent 需要检索文献、生成实验计划、执行脚本、汇总报告，天然是 open-ended 且长时链路。课程中的 ``evaluation is the key'' 在这类场景体现最明显：没有多维评估，团队会很快被 ``生成内容看起来不错'' 的假象误导。

\begin{warningbox}{研究助理 Agent 的幻觉风险更隐蔽}
这类系统经常出现 ``论证结构完整但证据链断裂'' 的输出。若评估只看语言流畅度或主观可读性，模型会被激励去写更像论文的文本，而非更真实的研究结论。
\end{warningbox}

\subsection{三条路径的共性工程清单}
\begin{enumerate}
    \item 必须先定义任务完成的机器可判定标准，再设计训练数据。
    \item 必须在评估中记录失败恢复路径，而不仅是最终成败。
    \item 必须把推理成本和人工复核成本一起纳入 ROI。
    \item 必须有最小权限原则和紧急回滚通道。
    \item 必须做版本化评估集，避免指标漂移导致误判。
\end{enumerate}

\begin{importantbox}{课程观点在实战中的映射}
不论是哪一类 Agent 项目，真正的交付单位都不再是 ``模型版本''，而是 ``模型 + 评估 + 系统 + 运行策略'' 的组合版本。只有组合版本才具有可复现业务价值。
\end{importantbox}

\subsection{本章小结}
三类项目虽然目标不同，但工程规律一致：先评估、再训练；先约束、再扩权；先可复现、再追求规模。课程中的方法论可以直接作为项目启动模板。


